\documentclass[10pt,a4paper]{amsart}
\usepackage[margin=19mm]{geometry}
\usepackage{amsmath,amssymb,mathtools}
\usepackage[scr=rsfs,cal=euler]{mathalfa}
\usepackage{xfrac}
\usepackage[unicode,hidelinks,bookmarks=false]{hyperref}
\newcommand{\ie}{i.e.\ }
\newcommand{\cf}{cf.\ }
\newcommand{\resp}{respectively\ }
\newcommand{\Subsection}{\subsection}
\providecommand{\nobreakdash}{\nobreak\hbox{-}}
\setlength{\emergencystretch}{2em}
\multlinegap0pt
\usepackage{amssymb}
\newtheorem{assumption}{Assumption}
\newtheorem{proposition}{Proposition}
\title{Global Well-Posedness for the Fourth-Order Nonlinear Schr\"{o}dinger Equation with Potential in the Energy-Critical Case}
\author{Hikaru Nakayama}
\date{Source: arXiv:2603.13892v1 (CC BY 4.0)}
\begin{document}
\maketitle

\noindent\emph{Source and licence.} Global Well-Posedness for the Fourth-Order Nonlinear Schr\"{o}dinger Equation with Potential in the Energy-Critical Case. Version arXiv:2603.13892v1 (CC BY 4.0), distributed by arXiv under the Creative Commons Attribution 4.0 International licence, http://creativecommons.org/licenses/by/4.0/, as recorded in the arXiv licence metadata for this exact version and shown on the abstract page https://arxiv.org/abs/2603.13892v1. Full licence notice: \texttt{licenses/arxiv-2603.13892-CC-BY-4.0.txt}.\par\smallskip

\noindent\textit{Editorial scope.} Counted unit: the article's proposition criterion (source0.tex lines 698--706). The article's complete original proof of it is delivered with it (source0.tex lines 708--774, one \texttt{proof} environment of 67 lines). No mathematics is written, repaired or re-derived: the statement and the proof are the article's own lines, in the article's own notation, and no retained line is substituted or re-flowed.  Retained uncounted as the source-local context the proof needs: lines 41--43; lines 83--90; lines 307--314; lines 323--325; lines 337--339; lines 493--498; lines 601--621.  Nothing was omitted from any retained span, so the package carries no omission record.  No retained line contains a citation, so the wrapper carries no bibliography and no bibliography entry.  No retained line contains a figure environment, an image-inclusion command or drawing code, so no image is delivered and the package declares no asset dependency.  Editorial formatting only, in addition: an AMS \texttt{amsart} preamble that restates the article's own declarations and macros used by the retained lines, namely the environments \texttt{assumption}, \texttt{proposition} on separate counters, together with the standard packages those lines need.  The retained environments print in this document as Assumption 1, Proposition 1 in order of appearance, whereas the article prints the same items with the numbering of its own full document.  The complete native source of the article as distributed in the arXiv e-print for this exact version is bundled unchanged as \texttt{sources/196356/source0.tex}.
\begin{equation}\label{eq:NL4S}
  i\partial_{t}u + \Delta^{2}u + Vu + \lambda|u|^{p-1}u = 0,\ \ x \in \mathbb{R}^{n}, \ \ t \in \mathbb{R}.
\end{equation}

\begin{assumption}\label{as:equ}
Let $V$ be a real-valued potential.
\begin{enumerate}
\item $V \geq 0$.
\item $V \in L^{\infty} \cap L^{\frac{n}{4}}$.
\item There exists a small $\delta_{n} > 0$ such that $\|V\|_{L^{\frac{n}{4},\infty}} \leq \delta_{n}$.
\end{enumerate}
\end{assumption}

\begin{assumption}\label{as:strichartz}
Let $V$ be a real-valued potential satisfying $V \in L^{\frac{n}{4}, \infty}$. Assume that for each $\ell = 0, 1, 2$, we have $(x\cdot \nabla)^{\ell}V \in L^{1}_{loc}$ and $|(x\cdot \nabla)^{\ell}V|^{1/2}(\Delta^{2} + 1)^{-1/2}$ is bounded on $L^{2}$. Furthermore, assume that for $u \in C_{0}^{\infty}$,
\begin{align}
  \langle (\Delta^{2} + V)u, u \rangle &\gtrsim \langle \Delta^{2}u, u \rangle, \\
  \langle (H_{1} + V_{1})u, u \rangle &\gtrsim \langle \Delta^{2}u, u \rangle, \\
  |\langle (H_{2} + V_{2})u, u \rangle| &\lesssim \langle (H_{1} + V_{1})u, u \rangle.
\end{align}
\end{assumption}

\begin{equation}\label{standard}
\|u\|_{L^{p_{1}}(I, L^{q_{1}})} \lesssim \|u_{0}\|_{L^{2}} + \|h\|_{L^{p_{2}^{\prime}}(I, L^{q_{2}^{\prime}})}.
\end{equation}

\begin{equation}\label{strichartz}
\|\Delta u\|_{L^{q}(I, L^{r})} \lesssim \|\Delta u_{0}\|_{L^{2}} + \|\nabla h\|_{L^{2}(I, L^{\frac{2n}{n + 2}})}.
\end{equation}

\begin{equation}\label{mass energy con}
\begin{aligned}
M(u(t)) &= M(u_{0}), \\
E(u(t)) &= E(u_{0}),
\end{aligned}
\end{equation}

\begin{proposition}\label{local existence}
Assume that $V$ satisfies Assumptions \ref{as:equ} and \ref{as:strichartz}. Let $n \geq 5$, and let $u_{0} \in H^{2}$ be arbitrary. For any interval $I = [0, T]$, there exists $\delta > 0$ such that, if
\begin{equation}\label{small}
\|e^{itH}u_{0}\|_{W(I)} < \delta,
\end{equation}
then there exists a unique solution $u\in C(I, H^{2})$ to \eqref{eq:NL4S} with initial data $u_{0}$, and moreover
$u\in M(I) \cap L^{\frac{2(n + 4)}{n}}(I \times \mathbb{R}^{n})$.
In addition, the mass and energy are conserved in the sense of \eqref{mass energy con}. Furthermore,
\begin{equation}\label{st}
\begin{array}{l}
\|u\|_{W(I)} < 2 \delta, \\
\|u\|_{M(I)} + \|u\|_{L^{\infty}(I, H^{2})} \leq C\left(\|u_{0}\|_{H^{2}} + \delta^{\frac{n + 4}{n - 4}}\right)
\end{array}
\end{equation}
holds for some constant $C > 0$ independent of $u_{0}$.
Moreover, depending on $\delta$, there exists $\delta_{0}$ such that for any $\delta_{1} \in (0, \delta_{0})$, if $\|u_{0} - v_{0}\|_{H^{2}} < \delta_{1}$ and $v$ denotes the solution to \eqref{eq:NL4S} with initial data $v_{0}$, then $v$ is defined on $I$ and, for any B-admissible pair $(q, r)$,
\begin{equation}\label{initconti}
\|u - v\|_{L^{q}(I, L^{r})} \leq C\delta_{1}
\end{equation}
holds, where $C > 0$ is independent of $u_{0}, v_{0}$.
\end{proposition}

\begin{proposition}\label{criterion}
Assume that $V$ satisfies Assumptions \ref{as:equ} and \ref{as:strichartz}. Let $n \geq 5$ and $p = 2^{\sharp} - 1$.
Let $u\in C([0, T), H^{2})$ be a solution to \eqref{eq:NL4S} satisfying $\|u\|_{Z([0, T])} < +\infty$.
Then there exists a constant $K = K(\|u_{0}\|_{H^{2}}, \|u\|_{Z([0, T])})$ such that
\begin{equation}\label{criterion:con:1}
\|u\|_{L^{\frac{2(n + 4)}{n}}([0, T], L^{\frac{2(n + 4)}{n}})} + \|u\|_{L^{\infty}([0, T], \dot{H^{2}})} + \|u\|_{M([0, T])} \leq K,
\end{equation}
and $u$ can be extended to a solution $\tilde{u}\in C([0, \tilde{T}), H^{2})$ of \eqref{eq:NL4S} for some $\tilde{T} > T$.
\end{proposition}

\begin{proof}
Take a small $\eta > 0$, and let $B = \|u\|_{Z([0, T])}$. For $x \geq 0$, let $[x]$ denote the largest integer not exceeding $x$.
Set
\[
N = \left[ (B/\eta)^{\frac{2(n+4)}{n-4}} \right] + 1.
\]
Then the interval $[0, T]$ can be decomposed into $N$ subintervals $I_{j}$, $j=1,\dots,N$, such that
\[
\|u\|_{Z(I_{j})} = \eta \quad (1 \leq j \leq N - 1), \qquad \|u\|_{Z(I_{N})} \leq \eta.
\]

Applying the Strichartz estimate \eqref{strichartz} on $I_{j} = [t_{j}, t_{j+1}]$, we obtain for any $t \in I_{j}$,
\begin{equation}\label{criterion:0}
\begin{aligned}
\|u\|_{M([t_{j}, t])}
&\leq C\|u(t_{j})\|_{\dot{H}^{2}} + C\|u\|_{Z(I_{j})}^{\frac{8}{n-4}}\|u\|_{M([t_{j}, t])} \\
&\leq C\|u(t_{j})\|_{\dot{H}^{2}} + C\eta^{\frac{8}{n-4}}\|u\|_{M([t_{j}, t])}.
\end{aligned}
\end{equation}
Also, applying \eqref{standard} on $I_{j}$ and using conservation of mass, we get
\begin{equation}\label{criterion:1}
\begin{aligned}
\|u\|_{L^{\frac{2(n+4)}{n}}([t_{j}, t], \ L^{\frac{2(n+4)}{n}})}
&\leq C\|u(t_{j})\|_{L^{2}} + C\|u\|_{Z([t_{j},\ t])}^{\frac{8}{n-4}}\|u\|_{L^{\frac{2(n+4)}{n}}([t_{j}, t], \ L^{\frac{2(n+4)}{n}})} \\
&\leq C\|u_{0}\|_{L^{2}} + C\eta^{\frac{8}{n-4}}\|u\|_{L^{\frac{2(n+4)}{n}}([t_{j}, t], \ L^{\frac{2(n+4)}{n}})}.
\end{aligned}
\end{equation}
If $\eta$ is chosen sufficiently small, \eqref{criterion:0} and \eqref{criterion:1} imply
\begin{align*}
\|u\|_{M(I_{j})} &\leq C\|u(t_{j})\|_{\dot{H}^{2}}, \\
\|u\|_{L^{\frac{2(n+4)}{n}}(I_{j}, \ L^{\frac{2(n+4)}{n}})} &\leq C\|u_{0}\|_{L^{2}}.
\end{align*}
Moreover, by \eqref{strichartz},
\[
\|u\|_{L^{\infty}(I_{j},\ \dot{H}^{2})} \leq C\|u(t_{j})\|_{\dot{H}^{2}},
\]
and in particular,
\[
\|u(t_{j+1})\|_{\dot{H}^{2}} \leq C\|u(t_{j})\|_{\dot{H}^{2}}.
\]
Thus, for every $j$,
\begin{equation}\label{criterion:2}
\begin{aligned}
&\|u\|_{L^{\frac{2(n+4)}{n}}([0, T], \ L^{\frac{2(n+4)}{n}})} \leq N^{\frac{n}{2(n+4)}}C\|u_{0}\|_{L^{2}}, \\
&\|u\|_{L^{\infty}([0,T],\ \dot{H}^{2})} \leq C^{N}\|u_{0}\|_{\dot{H}^{2}}, \\
&\|u\|_{M([0,T])} \leq C^{N}\|u_{0}\|_{\dot{H}^{2}}.
\end{aligned}
\end{equation}
Hence \eqref{criterion:con:1} follows from \eqref{criterion:2}.

Now take $t_{0} \in I_{N}$. By Duhamel's formula and the Strichartz estimate \eqref{strichartz}, for any $t$,
\begin{equation}\label{criterion:3}
\begin{aligned}
\|e^{i(t-t_{0})H}u(t_{0})\|_{W([t_{0}, T])}
&\leq \|u\|_{W([t_{0}, T])} + C\||u|^{\frac{8}{n-4}}u\|_{N([t_{0}, T])} \\
&\leq \|u\|_{W([t_{0}, T])} + C\|u\|_{W([t_{0}, T])}^{\frac{n+4}{n-4}}.
\end{aligned}
\end{equation}
Since the $W([t_{0}, T])$-norm of $u$ is finite, it can be made arbitrarily small by taking $t_{0}\to T$.
Therefore the left-hand side of \eqref{criterion:3} tends to $0$ as $t_{0}\to T$.
Hence, for any $\delta > 0$, there exist $t_{1} \in (0, T)$ and $T' > T$ such that $u(t_{1}) \in H^{2}$ and
\begin{equation}\label{criterion:4}
\|e^{i(t-t_{1})H}u(t_{1})\|_{W([t_{1}, T'])} \leq \delta.
\end{equation}
Then, by Proposition \ref{local existence}, there exists a solution $v \in C([t_{1}, T'], H^{2})$ of \eqref{eq:NL4S} with $v(t_{1}) = u(t_{1})$.
By uniqueness, $u = v$ on $[t_{1}, T)$, and thus $u$ can be extended to $[0, T']$.
\end{proof}

\end{document}
